\documentclass[10pt]{article}

\usepackage{amsmath}
%
\usepackage{amssymb}
%
\usepackage[width=7in,height=10in,top=0.75in,centering]{geometry}
%
\usepackage[dvipsnames]{xcolor}
%
\usepackage{hyperref}
%
%%%%%%%%%%%%%%%% HEADER %%%%%%%%%%%%%%%%%%%%%
\usepackage{fancyhdr}
\pagestyle{fancy}
%%
%\lfoot{Left footer}
%\cfoot{Center of footer}
%\rfoot{Right of footer}
%%
\lhead{Survey of Calculus I}
\chead{Knowledge Checkpoint 2 Reflection}
\rhead{\thepage}
%
\parindent0pt
\fboxsep=0.25cm
%
\newcommand{\Heading}[2]{\fbox{{\Large\textbf{#1}}}\hfill\fbox{\textbf{#2}}\par\vspace{0.1in}}
%
\newcommand{\Name}[0]{\textbf{Name}:\,\hrulefill\hrulefill\,\,}
%
\newcommand{\Date}[0]{\textbf{Date}:\,\hrulefill\,\,}
%
\newcommand{\Section}[0]{\textbf{Section}:\,\hrulefill\,\,}

%%%%%%%%%%%% BEGIN DOCUMENT %%%%%%%%%%%%%%%
%%%%%%%%%%%% BEGIN DOCUMENT %%%%%%%%%%%%%%%
%%%%%%%%%%%% BEGIN DOCUMENT %%%%%%%%%%%%%%%
%%%%%%%%%%%% BEGIN DOCUMENT %%%%%%%%%%%%%%%
%%%%%%%%%%%% BEGIN DOCUMENT %%%%%%%%%%%%%%%

\begin{document}

\thispagestyle{empty}

\Heading{MATH 1190/1210 Survey of Calculus I}{\begin{minipage}{0.25\textwidth}\begin{center} Knowledge Checkpoint 2\\Reflection \end{center}\end{minipage}}

\vspace{0.1in}

\Name{}\Date{}\Section{}

\hrulefill
%
\paragraph*{GOAL:} Struggle and making mistakes are healthy parts of growth as long as both are framed productively. The intention of this activity is to create an opportunity for you to re-frame struggle \& mistakes productively so you can grow as a learner.\\

{\bf That you are honest with yourself when you respond to these prompts is a critically important for your growth.} This can be difficult, so please just do your best.\\

\hrulefill\\

\begin{enumerate}

\item Read through the Knowledge Checkpoint 2 solutions on our Canvas site, and briefly review course materials. {\bf After doing so, which course experiences} (e.g. pre-class assignments, lessons, homework, practice problems) {\bf do you think would have best prepared someone for this assessment?}

    \vfill

\item Now read through your Knowledge Checkpoint 1 scores \& feedback. Did you improve your scores on learning targets reappearing on Knowledge Checkpoint 2? If so, then what did you do to promote your growth? If not, then what could you have done differently to improve?

\vfill

Regardless, please know that {\it making mistakes is \underline{not} a sign of weakness} and {\it mistakes do not reflect poorly on you as a person}. {\bf Mistakes are just opportunities to  grow!} :)

\newpage

\item {\bf What specific things did you do to prepare for Knowledge Checkpoint 2?}\\

\begin{minipage}[t][3.25in][t]{0.45\textwidth}
office hours \dotfill $\square$\\[0.1in]
MCLC \dotfill $\square$\\[0.1in]
P2L \dotfill $\square$\\[0.1in]
study group \dotfill $\square$\\[0.1in]
review pre-class videos \dotfill $\square$\\[0.1in]
``How to Study" videos \dotfill $\square$\\[0.1in]
spread out studying over at least a week \dotfill $\square$\\[0.1in]
cramming \dotfill $\square$\\[0.1in]
made a study schedule \dotfill $\square$\\[0.1in]
reread notes or textbook \dotfill $\square$\\[0.1in]
connect calculus to other courses \dotfill $\square$\\[0.1in]
\end{minipage}
\hfill
\begin{minipage}[t][3.25in][t]{.45\textwidth}
review WebAssign homework \dotfill $\square$\\[0.1in]
PAC \dotfill $\square$\\[0.1in]
lesson answer keys \dotfill $\square$\\[0.1in]
lesson extra practice problems \dotfill $\square$\\[0.1in]
learning target example problems \dotfill $\square$\\[0.1in]
learning target extra practice \dotfill $\square$\\[0.1in]
learning target concept checks \dotfill $\square$\\[0.1in]
multitasking while studying \dotfill $\square$\\[0.1in]
flashcards \dotfill $\square$\\[0.1in]
test yourself \dotfill $\square$\\[0.1in]
other (describe below) \dotfill $\square$\\[0.1in]
\end{minipage}
\hfill\,\\

In his third video, Dr. Chew tells us that successful learning isn't just a matter of developing more effective study skills: it involves overcoming ineffective or counter-productive skills that are highly practiced and automatic. For example, rereading notes and memorizing problems are generally counter-productive study skills but give students a sense that they're learning partly because these skills are easy. This can lead to the ``labor-in-vain" phenomenon in which students spend large amounts of time studying with essentially no learning gains.\\

Do you feel like you sometimes ``labor-in-vain" while studying for calculus? Which of your study practices do you think contribute to this feeling?

\vfill

%%%%%%%%%
\newpage%
%%%%%%%%%

    
\item In the article {\it\color{blue}\href{https://learninglab.psych.purdue.edu/downloads/2009/2009_Karpicke_Butler_Roediger.pdf}{Metacognitive Strategies in Student Learning...}}, the authors write: ``The testing effect refers to the finding that taking a test enhances long-term retention more than spending an equivalent amount of time repeatedly studying. There are clear and direct implications of the testing effect for student learning. One way for students to enhance their learning would be to practise recalling information while studying."\\

What are some new study practices you can adopt to help prevent you from experiencing ``labor-in-vain"? What are some ways you can modify your existing study practices to enhance their efficacy? 

    \vfill

\item I have a healthy, productive attitude about learning mathematics. (bubble one) \,\hfill $\bigcirc$ agree \hfill $\bigcirc$ disagree \hfill\,\\[0.1in]
%
Please write anything else you'd like to include in your assessment 2 exam wrapper.

    \vfill

\item The MATH 1190 \& 1210 instruction team wants you to know {\bf we believe in your ability to do great things with enough time and intentional effort}, and {\bf we want to support you in growing throughout this course}. We care, so don't hesitate to reach out if you have any questions or need any help!\\ 

{\bf Please acknowledge this message by signing your name below.}\\[0.5in]
    
\end{enumerate}

\end{document}